A rigorous investigation of structural aliasing requires that the claims be stated in a form that is both precise and empirically falsifiable.
Vague assertions that "patch-based models struggle with periodic signals" cannot guide principled experimental design or statistical evaluation.
This section therefore translates the theoretical results of Sections~\ref{sec:one}–\ref{sec:point3} into a set of concrete, testable hypotheses, each of which admits a clear null and alternative form suitable for Bayesian comparison in Section~\ref{sec:point7}.

\begin{description}
    \item[$H_1$ — Harmonic blind spots.]
    Forecasting error is significantly higher at signal frequencies that are integer multiples of $f_{\mathrm{patch}} = f_s / S$ than at off-harmonic control frequencies matched in amplitude and noise level.
    The null hypothesis $H_1^0$ states that error is identically distributed across all tested frequencies.

    \item[$H_2$ — Phase sensitivity.]
    At blind-spot frequencies, the model's prediction error varies systematically with the initial phase $\phi$ of the input sinusoid.
    Specifically, phase shifts that are not multiples of $2\pi$ relative to the patch grid are predicted to break the degeneracy and partially restore observability.
    The null hypothesis $H_2^0$ states that error is phase-invariant.

    \item[$H_3$ — Patch parameter dependence.]
    The location and severity of blind spots are determined by $P$, $S$, and their interaction $\mathrm{lcm}(P,S)$, as derived in Section~\ref{sec:two}.
    Increasing the overlap $P - S$ is predicted to reduce (but not eliminate) the aliasing effect at harmonic frequencies.
    The null hypothesis $H_3^0$ states that varying $P$ and $S$ (while holding $f_s$ fixed) does not alter the error profile across frequencies.
\end{description}

These three hypotheses form the experimental backbone of this work.
Additional exploratory hypotheses concerning the representational level (internal activation collapse) and the semantic level (ontology-class misclassification rate) are introduced in Section~\ref{sec:point5} and evaluated in Sections~\ref{sec:point6}–\ref{sec:point7}.
